\subsection{Mailman-style shared sums: arithmetic reuse meets scatter cost}
\label{sec:8-mailman}

The sign-code form of primary distance suggests reusing sums across
candidates that share bit patterns. The Mailman algorithm supplies the
general finite-alphabet matrix--vector idea \cite{mailman}.
The archived pilot tests GPU implementations on frozen Vogue, arXiv and
OpenAI-3072 candidate streams, at 128 probes and batches 8/1,000. Its
baseline reproduces the packed two-table scoring loop. Atomic bucket
accumulation, private histograms and group widths 4/6/8 are alternatives
within the experiment; the mathematical score and input codes remain fixed.

The fastest alternative in each of the six cells achieves only
0.187--0.581 of baseline scoring throughput (\cref{fig:8-negative}).
The phase profiles place most instrumented cycles in bucket accumulation.
Biased list-local patterns create repeated bucket destinations; private
histograms avoid atomics at the cost of larger state and shared-memory
read--modify--write work. Reducing group width and avoiding dominant buckets
improves particular variants, but none reaches the baseline. This supports
a stopping decision for these implementations and inputs. It does not
establish that all finite-alphabet reuse is ineffective: the winning
two-table baseline already reuses coordinate sums in a different direction.

The pilot writes primary scores rather than returning final neighbours.
Its FP32/FP64 comparisons expose near-cutoff numerical differences, not
final-pipeline equivalence. With measured scorer speed ratio $s$ and an
estimated replaceable latency fraction $f$, the optimistic replacement
model is $1/[(1-f)+f/s]$. Its predicted slowdown is a model, since integration
was stopped before an end-to-end run. An auxiliary lossless dictionary
study also stopped at structural statistics: its limited byte/work savings
did not justify a GPU implementation. Neither route contributes a measured
system point to the main frontiers.

\par\medskip\noindent\begin{minipage}{\linewidth}
\centering\captionsetup{hypcap=false}
\includegraphics[width=\linewidth]{ch8/negative-evidence.pdf}
\captionof{figure}{Two distinct evidence boundaries. Left: best measured
Mailman scorer / packed-table scorer speed on an RTX 5090, with Vogue,
arXiv and OpenAI-3072 at two batches and 128 probes; all ratios are below
one. Right: minimum survivor
counts meeting the original JHQ recall target on arXiv's 400-query FP64
offline test split. The candidate reduction is not a GPU speedup. Both
panels are regenerated from archived records by the thesis audit script.}
\label{fig:8-negative}
\end{minipage}\par\medskip

\subsection{Rank-oriented coding: a local gain is not a general replacement}
\label{sec:8-rank}

The rank-oriented study asks whether changing one-bit thresholds or scoring
can retain useful neighbours with fewer survivors at the same nominal code
budget. Its main comparisons replay exported candidates in FP64, with
400 calibration, 200 validation and 400 final-test queries. It examines
likelihood weighting, median thresholds, learned ranking losses,
linear-MMSE decoding and list-dependent thresholds. Different objectives
expose different failure modes: reducing reconstruction or inner-product
MSE can worsen the order of the few nearest candidates; changing primary
levels also changes the residual reference point.

On arXiv's test split, median thresholds reduce the minimum tested survivor
count meeting the original target from 875 to 375. That success does not
generalise to the other two workloads under the same strict criterion.
The shared-rotation Gaussian intuition also leaves real corpus bias intact:
rotating a mean vector preserves its norm. Centered coding can address some
of this mismatch, but its reconstructed-score ceiling must be evaluated
alongside primary ranking. The stopped study therefore supports a
workload-specific representation improvement, rather than a new generally
rank-optimal quantiser.

The later median-only GPU follow-up is separate evidence. It imports the
changed thresholds, levels and retrained residuals into the same search
pipeline, then compares independently swept budgets at measured recall
thresholds. At arXiv's original $\alpha=100$ target it yields ratios of
approximately 1.40 at batch 8 and 1.21 at batch 1,000. Vogue's tested median
frontier stops at 0.9630, below the original 0.9644 working point;
OpenAI-3072 is essentially tied. These results revise a blanket ``never
GPU-tested'' description, while preserving the decision against a universal
replacement. The follow-up uses its own index states, query pool and
three-round protocol; its rates are not inserted into E1.
